% Generated draft: requires mathematical and source review.
\documentclass[11pt]{article}
\usepackage[margin=25mm]{geometry}
\usepackage{fontspec}
\setmainfont{Latin Modern Roman}
\usepackage{amsmath,amssymb,mathtools,unicode-math}
\setmathfont{Latin Modern Math}
\usepackage{enumitem,needspace,xurl,hyperref}
\setlength{\parindent}{0pt}
\setlength{\parskip}{5pt}
\setlength{\emergencystretch}{3em}
\newcommand{\N}{\mathbb{N}}
\newcommand{\Z}{\mathbb{Z}}
\newcommand{\Q}{\mathbb{Q}}
\newcommand{\R}{\mathbb{R}}
\newcommand{\C}{\mathbb{C}}
\newcommand{\F}{\mathbb{F}}
\newcommand{\A}{\mathbb{A}}
\newcommand{\PP}{\mathbb{P}}
\newcommand{\E}{\mathbb{E}}
\newcommand{\eps}{\varepsilon}
\newcommand{\dd}{\,\mathrm{d}}
\begin{document}
\section*{UnsolvedMath ID: 30003324}
\textbf{Source problem code:} catalogue entry
\subsection{Definitions and mathematical statement}
Let \(n\geq 1\) be an integer. An integer-coefficient real exponential polynomial in \(n\) variables is a function \(f:\R^n\to\R\) given by a finite polynomial presentation
\[
P\in\Z[X_1,\ldots,X_n,Y_1,\ldots,Y_n],\qquad
f(x)=P(x_1,\ldots,x_n,e^{x_1},\ldots,e^{x_n}),
\]
where \(e^t\) is the usual real exponential function. Thus the notation \(\Z[x_1,\ldots,x_n,e^{x_1},\ldots,e^{x_n}]\) denotes these functions, supplied by such finite presentations. Write
\[
Z(f)=\{x\in\R^n:f(x)=0\},\qquad
\|x\|_2=\left(\sum_{i=1}^n x_i^2\right)^{1/2}.
\]
\textbf{First Root Conjecture.} There is a single algorithm which, on every input consisting of \(n\geq1\) and a finite polynomial presentation \(P\) as above, halts and returns a positive integer \(M(n,P)\) such that
\[
Z(f)\neq\varnothing
\quad\Longrightarrow\quad
\exists\alpha\in\R^n\;
\bigl(f(\alpha)=0\ \land\ \|\alpha\|_2<M(n,P)\bigr).
\]
The algorithm must halt also when \(Z(f)=\varnothing\); in that case the displayed implication imposes no further condition on its output. Only the existence of one zero within the bound is required, and the bound need not contain every zero.

\textit{Interpretive qualifications.} The supplied statement asks to effectively find \(\eta(n,f)\), without specifying its output type or encoding. The formulation above uses the usual uniform algorithmic interpretation with integer bounds and finite polynomial inputs. A positive integer bound is equivalent to a positive rational bound or a computably specified finite real bound that can effectively be enlarged to an integer. No canonical presentation of the function \(f\) is required. The source also leaves the norm unspecified; the Euclidean norm is used here as an explicit convention. Using the maximum norm yields an equivalent effective bounding problem, by
\[
\|x\|_\infty\leq\|x\|_2\leq\sqrt n\,\|x\|_\infty.
\]
These conventions require confirmation against the intended source formulation.
\subsection{Short English statement}
The First Root Conjecture asks for a uniform algorithm producing a finite radius from an integer-coefficient exponential polynomial. Whenever the polynomial has any real zero, at least one zero must lie strictly inside that radius. The supplied record labels the problem open; this draft does not independently verify that status.
\subsection{Sources}
\begin{itemize}
\item Local UnsolvedMath registry: permanent ID \(30003324\), title \emph{First Root Conjecture for Real Exponential Polynomials}, source problem code \texttt{catalogue entry}, in the \emph{Oberwolfach Reports Open Problems} collection. No individual proposer is supplied. The record gives proposed year 2016.
\item Supplied primary citation: \emph{Mini-Workshop: Surreal Numbers, Surreal Analysis, Hahn Fields and Derivations} (2017), \emph{Oberwolfach Reports}, pp.\ 3313--3372, DOI \texttt{10.4171/owr/2016/60}. Original OWR report: \url{https://doi.org/10.4171/owr/2016/60}. The supplied background describes Oberwolfach Reports as documenting mathematical workshops and their open-problem discussions.
\item Supplied literature source, labelled ``Krapp thesis'': \url{https://d-nb.info/1202012558/34}.
\item Supplied literature source, labelled ``Krapp on prime-model embeddings'': \url{https://arxiv.org/abs/2302.01609}.
\item Literature assessment dated \texttt{2026-08-21}: the supplied curation labels the problem open and reports that the First Root Conjecture remains tied to the unresolved effective decidability problem for real exponentiation, with later work recording partial progress but no proof. This is supplied dated curation, not an independently verified literature assessment. Its time-sensitive status and cited claims require independent verification.
\item The registry supplies statement status \texttt{corrected\_verified} and reports that the official report passage was checked and rewritten as a self-contained problem preserving its mathematical content. This is a supplied verification claim; no source visit or independent mathematical verification is asserted here.
\end{itemize}
\end{document}
